(1) 由球壳定理，球外壳层对内部点的引力净作用为零；半径 $r$ 内包围的质量是 $M(r)=Mr^3/R^3$。于是 $\mathbf a=-GM(r)\mathbf r/r^3=-GM\mathbf r/R^3$，在球体内部与位置成正比，等价于各向同性线性恢复加速度。
(2) 取穿心隧道的有符号坐标 $x$，运动方程为 $\ddot x+\omega^2x=0$，其中 $\omega^2=GM/R^3=g_s/R$，$g_s=GM/R^2$ 为表面重力。若 $x(0)=R,\dot x(0)=0$，解为 $x(t)=R\cos(\omega t)$，不存在假想的中心引力发散。
(3) 第一次到球心要求 $\cos(\omega t)=0$，故 $t_{\rm center}=\pi/(2\omega)$；到达对侧表面 $x=-R$ 的时间为 $t_{\rm other}=\pi/\omega$。回到原表面的完整周期 $T=2\pi/\omega=2\pi\sqrt{R^3/(GM)}$，与轨道半径 $R$ 的理想近地圆轨道周期形式相同。
(4) 非穿心直隧道上令 $s$ 自弦中点量起，位置矢量满足 $|\mathbf r|^2=b^2+s^2$。加速度沿隧道的分量为 $-GMs/R^3$，故仍满足 $\ddot s+\omega^2s=0$。球面交点满足 $s_{\max}=\sqrt{R^2-b^2}$，从一端静止释放时 $s=s_{\max}\cos\omega t$；频率与弦的偏移 $b$ 无关。
